problem:  
Let \((X,d,\mu)\) be a complete, separable, *doubling* metric measure space supporting a \(p\)-Poincaré inequality for some \(p>1\) (so that, by Cheeger’s theorem, \(X\) is a Lipschitz differentiability space).  
Assume that for \(\mu\)-a.e. \(x\in X\), every pointed measured Gromov–Hausdorff tangent \(\operatorname{Tan}(X,d,\mu,x)\) is isomorphic, as a metric measure space, to a *fixed* Carnot group \((G,d_G,\mu_G)\) with its Haar measure.  
Further, assume the tangents are *unique* up to isometry.  

**Open problem:**  
Does there exist a *measurable Carnot bundle structure* on \(X\), i.e. a Borel field of metric Lie groups \(\{G_x\}_{x\in X}\) with each \(G_x\cong G\), together with measurable identifications  
\[
\Phi_x:\operatorname{Lie}(G)\to\operatorname{Lie}(G_x),
\]
such that  
1. the group operation and dilation fields vary measurably in \(x\), and  
2. the corresponding pushforward of \(\mu\) and the tangent cone structure coincide with the measured Gromov–Hausdorff limits \(\operatorname{Tan}(X,d,\mu,x)\)?  

Equivalently, does almost-everywhere uniqueness and constancy of Carnot-group tangents force the infinitesimal geometry of \(X\) to assemble into a Borel measurable \(G\)-structure over \(X\)?  
If not, can one characterize the measurable obstructions (analogous to topological obstructions for principal \(G\)-bundles) arising from the automorphism group of \(G\)?

Why is it a "good" problem:  
This version articulates a deep, currently unresolved question at the heart of the intersection between metric measure geometry, sub-Riemannian analysis, and measurable groupoid theory. Removing the smooth manifold hypothesis exposes the true difficulty: in nonsmooth differentiability spaces, tangents may exist uniquely and be isomorphic to a fixed Carnot group, yet there is no known mechanism to glue these infinitesimal models into a coherent measurable field.  

The problem probes whether *synthetic constancy of metric-measure tangents* implies a measurable infinitesimal homogeneity—an analogue of the existence of tangent bundles in Riemannian geometry, but for *nonabelian, nonlinear* Carnot models. A positive resolution would extend Cheeger’s differentiable structure theory from Euclidean to Carnot settings, profoundly impacting the study of analysis and PDEs on metric spaces. A negative resolution would uncover new measurable invariants obstructing infinitesimal trivialization, potentially leading to a notion of “measurable curvature” in the sub-Riemannian category.  

The question is stated in simple geometric terms but reaches deeply into several modern domains—metric geometry, analysis, and group theory—making it a prototypical **“narrow entrance, profound depth”** research problem.